\scanpage{9}
Considérons maintenant l'immersion projective
\[
C\longrightarrow P(\mathcal E_3)
\]
définie par $L^{\otimes3}$ en vertu de prop. 1. Alors, $C$ devient un diviseur positif sur $P(\mathcal E_3)$ simple sur $S$, \inserted{de degré relatif 3,} il est donc défini par une \deleted{\unsure{élément|diviseur}} \inserted{section sur $S$} de l'espace des 3-formes
\[
\Phi_3(P(\mathcal E_3))\simeq P(\mathrm{Symm}_3(\check{\mathcal E}_3)).
\]
\deleted{De plus, comme}

\deleted{\uline{Supposons} que l'on se soit donné une trivialisation}

Considérons la section $s$ de $P(\mathcal E_3)$ et la \deleted{``droite'' \unsure{sur espace de dim 1|sous espace de dim 1}} \inserted{droite relative $D$} de $P(\mathcal E_3)$ définis comme $P(\mathcal E_3/\mathcal E_2)$ et $P(\mathcal E_3/\mathcal E_1)$, alors on constate que $s$ est une section de $C$, et que $D$ est une ``tangente d'inflexion'' ..... peu importe.

\figure{\begin{tikzpicture}[x=1cm,y=1cm]
\draw[line width=1.05pt] (0,0.12) -- (2.65,1.29);
\draw[line width=1.6pt,line cap=round] (.79,.47) -- (1.17,.64);
\draw[line width=.6pt] (.62,0) .. controls (.72,.60) and (1.20,.61) .. (1.53,.80) .. controls (1.98,1.00) and (2.27,1.28) .. (2.34,1.70);
\end{tikzpicture}}

Considérons une \deleted{trivialisation} \inserted{base de $\mathcal E_3$ compatible avec} \unsure{du clivage|des clivages} $0\subset\mathcal E_1\subset\mathcal E_2\subset\mathcal E_3$ de $\mathcal E_3$, \inserted{et commençant par $1_S\in\Gamma(S,\underline O_S)$,} [elle ``dépend en $(3+3-1)\inserted{{}=5}$ paramètres''], \deleted{par des coordonnées $\unsure{z,y,x|x,y,x}$} \inserted{i.e. une base $1,y,x$}

\[
\left(\begin{array}{ccc}
1&\unsure{\alpha|a}&\beta\\
\diagdown&\unsure{\mu|h}&\gamma\\
&&\unsure{\nu|v}
\end{array}\right)
\]

\deleted{Alors la forme $\varphi$ (définie à \deleted{\unsure{l'ordre}} \inserted{multiplication près par une} unité) prend la forme [où on remplace $z$ par $1$]
\[
a x^{\deleted{\unsure{3|2}}\unsure{2|3}}+b y^3+cxy+d y^2+ex+fy+g=0
\]
\unsure{$x^3y^2$|$y^3x^2$} où $y$ \unsure{désigne un suppl. de $\mathcal E_1$ dans $\mathcal E_2$} \inserted{\unsure{et $x$ de $\mathcal E_2$ dans $\mathcal E_3$}}}
% archive sheet number: 8
